\section{Interval normalization, continuation quotients, and graded splitting}
\label{sec:continuation-integration}

Report 14 contributes a different compression mechanism: change bases to
expose direct summands, then shorten certain entire summands when only a
final unknot verdict is required. The production implementation now includes
optional \code{barcode} and \code{fitting} recognition backends and exact-rank
interfaces. We retain the standard backend as default. The scalar splitter
is strengthened to respect recovered quantum shifts by default, allowing
its representation to retain the homogeneous coefficient invariant from
Section~\ref{sec:homogeneous-integration}.

\subsection{Exact decomposition of a pure component}
Let $m$ be one matching on $2p$ boundary endpoints. Its endomorphism ring is
\[
 R_m=\F[x_1,\ldots,x_p]/(x_1^2,\ldots,x_p^2).
\]
Every element $\theta$ with zero constant term satisfies $\theta^2=0$:
squares of nonconstant monomials vanish, and cross terms cancel in
characteristic two. Call a connected differential component pure if all
its objects are $m$ and every nonzero entry is one fixed nonzero
$\theta$ of this form. Singletons are also considered pure.

Write a pure component's degreewise maps as $\theta B_h$, where $B_h$
is binary. Crucially $B_{h+1}B_h$ need not vanish: it is the factor
$\theta^2$ that makes the actual differential square to zero. The $B_h$
form a representation of a finite linearly oriented quiver, whose interval
decomposition is a simultaneous scalar change of basis. Define
\[
 \rho(a,b)=\operatorname{rank}(B_{b-1}\cdots B_a),\qquad
 \rho(a,a)=\dim V_a.
\]
The number of intervals supported on $[a,b]$ is
\[
 \rho(a,b)-\rho(a-1,b)-\rho(a,b+1)+\rho(a-1,b+1),
\]
with unavailable ranks interpreted as zero. Indeed $\rho(a,b)$ counts
intervals covering $[a,b]$, and this finite difference selects exactly
the intervals born at $a$ and ending at $b$. Successive image-space bases
compute these ranks without materializing all matrix products.

Reinserting $\theta$ gives a chain isomorphism to a direct sum of intervals
$I_l(m,\theta)$, each with $l$ consecutive copies of $m$ and maps
$\theta$. The implementation reconstructs intervals in degree order and
then shares identical templates using the existing nonnegative
multiplicity polynomials. This decomposition preserves full raw homology
and every later additive continuation. It is applied to whole differential
components, never to a submatrix that still has unaccounted attachments.

In exact mode, lengths and birth shifts are retained. In decision mode,
birth shifts may be discarded and multiplicities saturated at three.
These are two distinct interfaces. In particular, scalar coefficients
remain in $\F$, whereas multiplicities are counts and must not be reduced
modulo two.

\subsection{Why length two is enough for the decision observation}
Put $A=\F[\epsilon]/(\epsilon^2)$ and let $J_l$ have $l$ free rank-one
terms with differential $\epsilon$. For an actual remaining crossing
sequence and ordinary closure $F$, put $D=F(m)$ and let $\epsilon$ act
by $F(\theta)$. This action commutes with the suffix differential and
squares to zero. In the strict raw complex, additive planar gluing gives
\[
 F(I_l(m,\theta))\cong J_l\otimes_A D.
\]
This module structure, not linearity alone, justifies the replacement.

For a bounded complex $D$ of finite-dimensional $A$-modules, the
finite dual-number derived classification gives
\begin{equation}\label{eq:integrated-continuation-rank}
 r_l(D)=\dim H(J_l\otimes_A D)
       =a l+2\sum_{j\ge1}b_j\min(l,j),
\end{equation}
where $a,b_j$ are nonnegative and finitely supported. Here $a$ counts
shifted simple modules and $b_j$ counts shifted free intervals $J_j$.
Report 14 supplies an elementary proof: take a bounded-above degreewise
finite free resolution and cancel its units downward. The remaining
matrices are $\epsilon B_h$. In an exact far-left tail,
$2r_h-\operatorname{rank}B_{h-1}-\operatorname{rank}B_h=0$ forces
constant rank and invertible $B_h$. Choose identity bases on that tail,
cut it once, and interval-decompose the finite remainder. Intervals born
after the cut stay finite; those born at the cut extend left and represent
shifted simple modules. Thus only finitely many summands are needed.

For completeness, the tensor calculation behind
\eqref{eq:integrated-continuation-rank} is
\[
 \dim H(J_l\otimes_A J_j)=2\min(l,j),\qquad
 \dim H(J_l\otimes_A\F)=l.
\]
The first differential is $\epsilon(u+v)$ on
$A\otimes\F[u,v]/(u^l,v^j)$. Its homology dimension is twice the
cokernel dimension of multiplication by $u+v$, namely $2\min(l,j)$.
The second differential vanishes. Since $J_l$ is bounded and free,
tensoring preserves the quasi-isomorphisms used above.

A nonzero radical endomorphism requires a nonempty matching. Every complete
smoothing of an actual suffix glued to it contains at least one circle,
so its unreduced chain-space dimension is even. Homology has the same
parity, implying $r_1(D)=a+2\sum b_j$ is even and therefore $a$ is even.
For $l\ge2$, if $a>0$ or any $b_j>0$ with $j\ge2$, then already
$r_2(D)\ge4$. Otherwise $r_l(D)=2b_1$ is constant. Consequently
\begin{equation}\label{eq:integrated-cap-two}
 \min(3,r_l(D))=\min(3,r_{\min(l,2)}(D)).
\end{equation}
Length one cannot replace length two uniformly: close a two-arc matching
to one circle and take $\theta=x_1+x_2$. The two dots become the same
circle operator and cancel, so the interval's closed rank is $2l$.
For arbitrary algebraic suffixes without the evenness premise, the simple
module gives $r_l=l$, and cutoff two fails. The optional algebraic cutoff
three remains valid in that broader setting.

Shortening preserves an observation, not the homotopy type or the exact
homology. Repeated replacements are sound because every later continuation
is also a continuation before the replacement. Direct sums and copies
preserve capped-rank equality, and saturation commutes with addition and
multiplication of nonnegative counts. At final closure, elimination leaves
isolated generators; their saturated total is exactly the original rank
capped at three. Rank two certifies the unknot, and a capped value three
means at least three, hence at least four for a classical knot.
The code labels shortened-model object statistics accordingly and never
exports their degree counts as exact knot homology.

\subsection{Certified scalar splitting with recovered grading}
Purity need not be visible in the current basis. The optional Fitting
search looks for scalar endomorphisms $Q$ satisfying $Qd=dQ$, solving the
binary equations coefficient by coefficient. The archive permits variables
between objects with equal matching and homological degree. We additionally
require equality of recovered quantum shift by default.

For a nonzero homogeneous entry $u\to v$, let $d_{uv}$ be its dot degree
and $c_{uv}$ the overlay-circle count. The required relation is
\[
 q_v-q_u=p-c_{uv}+2d_{uv}.
\]
Propagating these differences along the undirected differential graph
recovers all $q_v$ up to one additive constant per component. The code
checks homogeneity of every inspected entry and rejects inconsistent cycle
sums. Scalar variables exist only within blocks $(m,h,q)$. Both solving
and basis-change verification use these blocks. Cache keys include the
relative shifts, so a cached ungraded search cannot silently stand in for
a graded search or a geometrically different shift constraint.

For a candidate $Q$, take a power $Q^N$ with $N$ a power of two at least
the largest block dimension. Then
$\ker Q^N\oplus\operatorname{im}Q^N$ is the stable Fitting decomposition.
Because $Q$ commutes with $d$, these spaces form subcomplexes.
The code builds bases for both parts, checks invertibility, verifies the
commutator independently, transforms every differential coefficient, and
requires every entry across the proposed split to vanish. It accepts only
nontrivial splits. Optional witnesses record the original and transformed
matrices, matching and degree labels, relative quantum shifts, and the
change of basis. A separate dense test verifier checks both directions of
the basis transformation and rejects corrupted witnesses.

The search is deliberately incomplete: defaults limit each considered
component to 48 objects, the scalar system to 1024 variables, accepted
splits to 16 per crossing, and candidates to a bounded deterministic set.
A local size or search limit leaves the exact component intact. Accepted
splits are followed by interval normalization and sharing. Exhausting the
global scanner budget gives \code{UNKNOWN}. The global time allowance
includes order selection, and retained-object ceilings are checked before
crossing allocation. The low-level \code{preserve\_grading=False} option
retains the archive's exact ungraded search for controlled comparisons;
it does not justify the homogeneous bounds below.

For a connected pure component with homogeneous $\theta$ of dot degree
$d$, $q_v-2d h_v$ is constant throughout the component. Thus all objects
in a given homological layer have the same recovered quantum shift, and
its interval basis changes are automatically graded. Sharing identical
connected templates preserves the existence of a graded lift up to
componentwise shifts. Shortening keeps a graded interval representative,
although it does not preserve the original bigraded homology. These facts
justify using homogeneous coefficient counts after our restricted scalar
splits; no claim to output the erased absolute quantum grading is made.

\subsection{A sharper checkpoint bound}
A pure checkpoint is a stage where, after the implemented splitting and
normalization, every component is pure. Such a stage is recognizable; a
small maximum gap between such stages is an additional family hypothesis.
Let $w=2p$ bound the frontier throughout the chosen order, and let $g$
bound the number of crossings between successive pure checkpoints,
including the initial and final stages.

Without grading, there are $2^{2^p-1}-1$ possible nonzero radical
coefficients on a fixed matching. With our homogeneous invariant there are
at most
\[
 S(p)=\sum_{d=1}^p\left(2^{\binom pd}-1\right)
 \le p\,2^{\binom p{\lfloor p/2\rfloor}}\quad(p\ge1),
 \qquad S(0)=0.
\]
At a checkpoint, a shortened interval type is determined by its matching
and coefficient; singletons need no coefficient. With
$M(w)=(w-1)!!$ for $w>0$ and $M(0)=1$, there are at most
\[
 M(w)(1+S(p))\ \text{templates},\qquad
 M(w)(1+2S(p))\ \text{stored objects}.
\]
No common boundary disk is assumed. This bound is independent of the
original connected component size, the number of hidden copies, and the
original interval lengths. Exact mode retains lengths through $n+1$ and
pays only an additional polynomial factor in $n$.

One crossing makes at most eight delooped children per object. Hence an
impure stretch costs at most a factor $8^g$ in stored objects. Gaussian
elimination, interval ranks, the bounded scalar search, and coefficient
operations have polynomial cost in this physical size and $2^{O(w)}$.
Multiplicity polynomials have polynomially many degree positions and
$O(n)$-bit coefficients. The resulting conservative bound is
\begin{equation}\label{eq:graded-checkpoint-cost}
 T(n;w,g)\le n^{O(1)}
 \exp\left(O\left(g+w\log(w+1)+
                  \binom{w/2}{\lfloor w/4\rfloor}\right)\right).
\end{equation}
This sharpens the report's $2^{w/2}$ term by incorporating the recovered
grading. It gives a complete quasi-polynomial recognizer on families with
constructible suitable orders for which the displayed exponent is
$O(\log^2(n+2))$. It does not prove that arbitrary knot diagrams have such
orders or frequent pure checkpoints. A single gap $g=n$ leaves an
exponential bound. A normalization size cap that skips a nontrivial
component prevents that stage from counting as a pure checkpoint in the
implemented statistics.

\subsection{Validation and practical scope}
All 267 production tests pass. The imported interval tests independently
change quiver bases, reconstruct interval multiplicities, and compare dense
suffix closures. The scalar tests replay witnesses, check corruption, and
compare 80 seeded knot braids with exact full ranks and capped decisions;
another 60 seeded braid cases exercise interval normalization. New production
regressions cover grading constraints, inconsistent cycles, cache separation,
order-selection deadlines, input validation, CLI incompatibilities, and
recognition fallback after an agreeing window query. Recovered shifts also
agree up to componentwise constants with independently propagated shadow
labels on Conway, Kinoshita--Terasaka, and the hard unknot.

Report 14's independent continuation verifier imports no production code.
It enumerates 861 finite complexes of dual-number modules, including
nonfree modules and unit differentials, and passes 10\,206 threshold
comparisons through interval length six. This supports the implementation
and finite algebra examples; the all-suffix statement rests on the proof.
Its output and the full integrated test log are retained under
\path{synthesis/data/}.

\begin{center}
\begin{tabular}{@{}lrrrr@{}}
\toprule
Input & Interval exact & Fitting exact & Fitting decision & A/A \\
\midrule
conway & 0.821 & 0.460 & 0.469 & 1.063 \\
kinoshita\_terasaka & 0.841 & 0.487 & 0.481 & 0.992 \\
hard\_unknot\_8 & 0.819 & 0.777 & 0.795 & 1.028 \\
stress\_braid5\_36 & 0.747 & 0.776 & 0.766 & 0.988 \\
torus\_3\_5 & 0.901 & 0.889 & 0.871 & 1.003 \\
\bottomrule
\end{tabular}
\end{center}
Ratios are median paired full-scanner/new times; values below one are slower.


These are seven shuffled paired raw-backend rounds on CPython 3.13.14,
including order selection and excluding imports and diagram construction.
Every exact result is checked degree by degree against the full scanner;
every decision result is checked against its rank capped at three. A/A
compares identical full-scanner arms. Witness recording is performed
separately, outside the timed arms. Earlier certificates are bypassed, so
these measurements describe backend costs rather than end-to-end default
recognition. Default front-end methods and all inputs are recorded.

The data include overhead and unsuccessful normalization attempts. None
of these five actual diagrams produces an eligible nonsingleton pure
component, so they demonstrate no real-diagram interval-shortening gain.
The synthetic ladder regression does exhibit the intended structural
saving: 31 hidden copies of a seven-term interval require 217 objects under
literal sharing, seven after exact interval normalization, and two after
decision shortening. This is a graded algebraic fixture, not a claim that
it appears in the knot corpus. Scalar splitting produces two verified splits on both Conway and
Kinoshita--Terasaka, but their existence alone does not imply faster recognition.
On this corpus the new modes are slower than the full baseline. The quantum
restriction also costs time: the median ungraded/graded Fitting ratio is
about $0.85$ on those two diagrams and near one on the other three.
The measured tradeoffs support keeping both backends optional.

Report 14's bounded crossing-order optimizer is a separate proposal and
is not part of this integration. Its objective is a frontier proxy rather
than elapsed scanner time. The general quasi-polynomial recognition goal
still requires a structural theorem or additional machinery beyond these
conditional compression bounds.
